% GENERATED FILE. Edit canonical lessons, then run npm run generate:books.
% canonical-source-hash: fnv1a64:59c3bb35664f9a2a
% canonical-lessons: BN-C190-borong, BN-C190-hoyto, BN-C190-nishchoyi, BN-C190-durbhagyoboshoto, BN-C190-pholaphol

\chapter{Rather, Perhaps, Certainly, Unfortunately, Result}
\label{ch:bn-rather-perhaps-certainly-unfortunately-result}
\hlchaptermodality{190}

\begin{chapteropening}
\textbf{By the end of this chapter:} \emph{I can say rather, perhaps, certainly, unfortunately and a result.}

Complete the last lesson of chapter 190: I can say rather, perhaps, certainly, unfortunately and a result.
\end{chapteropening}

\section[\texorpdfstring{\bn{বরং} (bôrông)}{বরং (bôrông)}]{\bn{বরং} (bôrông) — rather}

\label{lesson:BN-C190-borong}



\begin{quote}
Before the new one: say the Bengali for a ravine, then the Bengali for a waterfall.
\end{quote}

\subsection*{You'll want to know: \bn{বরং}}
\textbf{\bn{বরং}} — \emph{bôrông} — \textquotedblleft{}rather\textquotedblright{}.

\begin{cousinweb}
\textbf{\bn{বরং}} (\emph{bôrông}) is Sanskrit \emph{varam}, \textquotedblleft{}better than\textquotedblright{}, taken over whole as a learned loan, a tatsama (CDIAL 11308). It is the adverb of \emph{vara}, \textquotedblleft{}choice, boon\textquotedblright{}, whose inherited form \emph{bar}, \textquotedblleft{}boon\textquotedblright{}, CDIAL lists.
\end{cousinweb}

\begin{grammarlens}[title={What you've built}]
One more word for this chapter.
\end{grammarlens}

\subsection*{Your turn}
\begin{itemize}
\raggedright
  \item \emph{Say it:} \emph{bôrông}
  \item \emph{Say it:} \emph{bôrông}, once more
  \item \emph{Say it:} say \emph{bôrông}
  \item \emph{Recall:} say the Bengali for a ravine, then the Bengali for a waterfall, then say \emph{bôrông} again
\end{itemize}

\begin{tcolorbox}[breakable,colback=teal!4,colframe=teal!35!black,title={Before you move on}]
What does \bn{বরং} mean? (\textquotedblleft{}Rather\textquotedblright{}.) Say it once more.
\end{tcolorbox}
\section[\texorpdfstring{\bn{হয়তো} (hôyto)}{হয়তো (hôyto)}]{\bn{হয়তো} (hôyto) — perhaps}

\label{lesson:BN-C190-hoyto}



\begin{quote}
Before the new one: say the Bengali for a waterfall, then the Bengali for rather.
\end{quote}

\subsection*{You'll want to know: \bn{হয়তো}}
\textbf{\bn{হয়তো}} — \emph{hôyto} — \textquotedblleft{}perhaps\textquotedblright{}.

\begin{cousinweb}
\textbf{\bn{হয়তো}} (\emph{hôyto}) is \emph{hôy}, \textquotedblleft{}is\textquotedblright{}, a form of \emph{hôoyā}, \textquotedblleft{}to be\textquotedblright{}, from Sanskrit \emph{bhavati}, \textquotedblleft{}becomes, is\textquotedblright{} (CDIAL 9416), + the particle \emph{-to}, \textquotedblleft{}indeed\textquotedblright{}: \textquotedblleft{}it may be so\textquotedblright{}. CDIAL's addenda move \emph{hôy} there from \emph{asti}, \textquotedblleft{}is\textquotedblright{} (CDIAL 977).
\end{cousinweb}

\begin{grammarlens}[title={What you've built}]
One more word for this chapter.
\end{grammarlens}

\subsection*{Your turn}
\begin{itemize}
\raggedright
  \item \emph{Say it:} \emph{hôyto}
  \item \emph{Say it:} \emph{hôyto}, once more
  \item \emph{Say it:} say \emph{hôyto}
  \item \emph{Recall:} say the Bengali for a waterfall, then the Bengali for rather, then say \emph{hôyto} again
\end{itemize}

\begin{tcolorbox}[breakable,colback=teal!4,colframe=teal!35!black,title={Before you move on}]
What does \bn{হয়তো} mean? (\textquotedblleft{}Perhaps\textquotedblright{}.) Say it once more.
\end{tcolorbox}
\section[\texorpdfstring{\bn{নিশ্চয়ই} (nishchôyi)}{নিশ্চয়ই (nishchôyi)}]{\bn{নিশ্চয়ই} (nishchôyi) — certainly}

\label{lesson:BN-C190-nishchoyi}



\begin{quote}
Before the new one: say the Bengali for rather, then the Bengali for perhaps.
\end{quote}

\subsection*{You'll want to know: \bn{নিশ্চয়ই}}
\textbf{\bn{নিশ্চয়ই}} — \emph{nishchôyi} — \textquotedblleft{}certainly\textquotedblright{}.

\begin{cousinweb}
\textbf{\bn{নিশ্চয়ই}} (\emph{nishchôyi}) is Sanskrit \emph{niścaya}, \textquotedblleft{}certainty\textquotedblright{}, taken over whole, + the stressing ending \emph{-i}, \textquotedblleft{}indeed\textquotedblright{}.
\end{cousinweb}

\begin{grammarlens}[title={What you've built}]
One more word for this chapter.
\end{grammarlens}

\subsection*{Your turn}
\begin{itemize}
\raggedright
  \item \emph{Say it:} \emph{nishchôyi}
  \item \emph{Say it:} \emph{nishchôyi}, once more
  \item \emph{Say it:} say \emph{nishchôyi}
  \item \emph{Recall:} say the Bengali for rather, then the Bengali for perhaps, then say \emph{nishchôyi} again
\end{itemize}

\begin{tcolorbox}[breakable,colback=teal!4,colframe=teal!35!black,title={Before you move on}]
What does \bn{নিশ্চয়ই} mean? (\textquotedblleft{}Certainly\textquotedblright{}.) Say it once more.
\end{tcolorbox}
\section[\texorpdfstring{\bn{দুর্ভাগ্যবশত} (durbhāgyôbôshôtô)}{দুর্ভাগ্যবশত (durbhāgyôbôshôtô)}]{\bn{দুর্ভাগ্যবশত} (durbhāgyôbôshôtô) — unfortunately}

\label{lesson:BN-C190-durbhagyoboshoto}



\begin{quote}
Before the new one: say the Bengali for perhaps, then the Bengali for certainly.
\end{quote}

\subsection*{You'll want to know: \bn{দুর্ভাগ্যবশত}}
\textbf{\bn{দুর্ভাগ্যবশত}} — \emph{durbhāgyôbôshôtô} — \textquotedblleft{}unfortunately\textquotedblright{}.

\begin{cousinweb}
\textbf{\bn{দুর্ভাগ্যবশত}} (\emph{durbhāgyôbôshôtô}) is built from Sanskrit \emph{durbhāgya}, \textquotedblleft{}ill luck\textquotedblright{}, and \emph{vaśataḥ}, \textquotedblleft{}by force of\textquotedblright{}, both taken over whole: \textquotedblleft{}by force of ill luck\textquotedblright{}.
\end{cousinweb}

\begin{grammarlens}[title={What you've built}]
One more word for this chapter.
\end{grammarlens}

\subsection*{Your turn}
\begin{itemize}
\raggedright
  \item \emph{Say it:} \emph{durbhāgyôbôshôtô}
  \item \emph{Say it:} \emph{durbhāgyôbôshôtô}, once more
  \item \emph{Say it:} say \emph{durbhāgyôbôshôtô}
  \item \emph{Recall:} say the Bengali for perhaps, then the Bengali for certainly, then say \emph{durbhāgyôbôshôtô} again
\end{itemize}

\begin{tcolorbox}[breakable,colback=teal!4,colframe=teal!35!black,title={Before you move on}]
What does \bn{দুর্ভাগ্যবশত} mean? (\textquotedblleft{}Unfortunately\textquotedblright{}.) Say it once more.
\end{tcolorbox}
\section[\texorpdfstring{\bn{ফলাফল} (phôlāphôl)}{ফলাফল (phôlāphôl)}]{\bn{ফলাফল} (phôlāphôl) — a result}

\label{lesson:BN-C190-pholaphol}



\begin{quote}
Before the new one: say the Bengali for certainly, then the Bengali for unfortunately.
\end{quote}

\subsection*{You'll want to know: \bn{ফলাফল}}
\textbf{\bn{ফলাফল}} — \emph{phôlāphôl} — \textquotedblleft{}a result\textquotedblright{}.

\begin{cousinweb}
\textbf{\bn{ফলাফল}} (\emph{phôlāphôl}) is Sanskrit \emph{phala}, \textquotedblleft{}fruit, result\textquotedblright{}, + \emph{aphala}, \textquotedblleft{}no result\textquotedblright{}, both taken over whole: \textquotedblleft{}result or no result\textquotedblright{}. CDIAL lists the inherited Bengali \emph{phal}, \textquotedblleft{}fruit\textquotedblright{} (CDIAL 9051).
\end{cousinweb}

\begin{grammarlens}[title={What you've built}]
One more word for this chapter.
\end{grammarlens}

\subsection*{Your turn}
\begin{itemize}
\raggedright
  \item \emph{Say it:} \emph{phôlāphôl}
  \item \emph{Say it:} \emph{phôlāphôl}, once more
  \item \emph{Say it:} say \emph{phôlāphôl}
  \item \emph{Recall:} say the Bengali for certainly, then the Bengali for unfortunately, then say \emph{phôlāphôl} again
\end{itemize}

\begin{tcolorbox}[breakable,colback=teal!4,colframe=teal!35!black,title={Before you move on}]
What does \bn{ফলাফল} mean? (\textquotedblleft{}A result\textquotedblright{}.) Say it once more.
\end{tcolorbox}
